\section{引言}

\subsection{背景：能力增长与可靠性缺口的错位}

过去数年间，大语言模型（Large Language Model, LLM）在数学推理任务上的表现提升显著。以思维链（Chain-of-Thought, CoT）提示为代表的推理激发技术\citep{wei2022chain,kojima2022large}，把模型的输出从“直接给答案”变为“先写推理再给答案”，并在 GSM8K\citep{cobbe2021training}、MATH\citep{hendrycks2021measuring} 等基准上带来了可观的准确率增益。随后的工作沿两条路径继续推进：一是把自然语言推理与外部工具耦合，让部分计算落到可执行程序上\citep{chen2022program,gao2022pal,schick2023toolformer}；二是把数学陈述翻译进形式系统，借助证明助手完成机器可检查的证明\citep{moura2021lean,mathlib2020lean,jiang2022draft,yang2023leandojo,xin2024deepseek}。

然而，能力增长与可靠性之间并不存在自动的同步关系。三个反复被观察到的现象值得注意。

\textbf{其一，答案正确不等于推理正确。}基准分数是对最终答案的匹配，而模型完全可能以错误的中间推理得到正确的终点。对同一批错误案例的重新标注显示，模型在“答案错”与“推理错”两个维度上的分布并不重合。

\textbf{其二，模型对自身推理的批评并不可靠。}针对“让模型自我纠错以提升推理”的流行做法，有工作系统地检验后发现：在缺乏外部反馈的情况下，自我纠错常常把已经正确的答案改成错误答案，整体表现不升反降\citep{huang2023large}。

\textbf{其三，性能对表面扰动敏感。}GSM-Symbolic 通过仅改变姓名与数值而对问题结构保持不变的方式扰动 GSM8K，模型准确率出现大幅下降，且随着扰动量增加继续下滑\citep{mirzadeh2024gsm}。这说明相当一部分“推理”更接近对训练分布中表面模式的复现，而非对结构的把握。

这三个现象指向同一个结构性缺口：\textbf{自然语言形式的推理过程，本身不可被机械验证}。验证者只能看到一串由模型生成的字符串，并用自己的语义理解去判断它对不对——而验证者与生成者共享同一套易错的语义判断，误差因此无法被独立检出。

\subsection{两条既有路线及其遗留问题}

围绕这一缺口，已有两条主要路线。

\textbf{路线一：过程级验证（process-level verification）。}不只看最终答案，而是训练一个验证器对推理的每一步打分\citep{lightman2023let,uesato2022solving}。这一路线在工程上有效，但存在一个根本的性质问题：验证器本身是学习得到的模型，它的判断同样可能出错，且其错误模式与生成模型高度相关。也就是说，验证器把“不可验证的自然语言判断”转移到了另一个不可验证的自然语言判断上，并未引入外部性。

\textbf{路线二：形式化验证（formal verification）。}把数学陈述与证明翻译进 Lean\citep{moura2021lean,mathlib2020lean}、Coq\citep{bertot2004interactive} 等证明助手，由内核做机器检查——这是目前可靠性最强的方案，其可靠性来自内核而非任何模型的判断\citep{wadler2015propositions}。这一路线面临两类困难：一是自动形式化（autoformalization）能力仍不足，自然语言到形式陈述的翻译本身是瓶颈\citep{szegedy2020promising,wu2022autoformalization}；二是即使翻译成功，形式系统的可覆盖范围也有限，大量数学论证（近似、估计、尺度分析、启发式构造）并不天然落在形式化边界之内。

两条路线因此各自留下了未被处理的空间：路线一缺乏外部性，路线二缺乏覆盖性。而在这两者之间，还存在一个既未被明确命名、也未被系统地作为研究对象处理的中间环节。

\subsection{本文的问题、主张与贡献}

\textbf{研究问题。}在不要求全流程形式化的前提下，如何让自然语言数学推理获得可检查、可复现的可靠性保障？

\textbf{核心主张。}本文主张：\emph{可靠推理 = 生成 + 语义规约 + 验证}，其中「\textbf{语义规约层}」把不可验证的自然语言推理降为可验证对象的\emph{唯一中间环节}，因而也是当前 AI4Math 系统的主要失效点与投入产出比最高的改进点。所谓语义规约，是指从自然语言推理出发，抽取并显式化其中的变量、假设、目标状态与状态变换，形成一份\textbf{类型化中间表示}（Typed Intermediate Representation, Typed IR）的过程。

\textbf{本文贡献。}本文为分析型与框架型研究，不含新训练或新微调，贡献如下：

\begin{enumerate}[leftmargin=1.6em,itemsep=2pt]
  \item \textbf{三层可靠性模型。}给出生成层（L1）、语义规约层（L2）、验证层（L3）的定义、输入输出与可验证性归属，并给出三条可检查的推论（第 \ref{sec:model} 节）。
  \item \textbf{失效模式分类（C-1）。}把 L2 的失败从“模型不够强”细化为四类可诊断模式：歧义未消解、类型不完整、不可逆、语义漂移（第 \ref{sec:failure} 节）。
  \item \textbf{验证深度梯度 V0--V3（C-2）与三项指标（C-3）。}给出按成本选择验证强度的阶梯，以及可规约率、规约保真度、反例检出率的定义式（第 \ref{sec:gradient} 节）。
  \item \textbf{实践者自检清单（C-4）。}一份十条清单，逐条对应上述失效模式（第 \ref{sec:checklist} 节）。
\end{enumerate}

\textbf{结论边界（先行声明）。}为避免结论超出证据，本文明确不做以下五件事：不训练或微调模型；不声称四类失效模式是完备的；不给出大样本实证（第三项贡献只给定义与最小可用示例）；不主张取代形式化路线，而定位为\emph{降低进入 L3 的门槛}；不声称适用于全部数学，范围限定在“离散、可类型化、步骤可逐步状态化”的推理类型。第 \ref{sec:limits} 节将逐条展开并复核正文是否越界。
